% !PRINTED 117
\weberpage{117}
%<<CONT
cyclic ($S = 1, \gamma, \gamma^2, \gamma^3, \gamma^4, \gamma^5$). There remains, then, only the assumption that $\gamma$ is of degree $2$, so that
\weq{2}{\gamma^2 = 1}
holds.
Since, however, we can replace $\gamma$ by $\gamma\alpha$ or $\gamma\alpha^2$, these too must be of degree $2$, and it follows that:
\weq{3}{\gamma\alpha = \alpha^2\gamma, \quad \gamma\alpha^2 = \alpha\gamma,}
with whose help one can always reduce every compositum of arbitrarily many powers of $\alpha$ and $\gamma$ to one and only one of the elements $S$. And if one now denotes
\[
1, \alpha, \alpha^2, \gamma, \gamma\alpha, \gamma\alpha^2
\]
by
\[
1, \alpha, \beta, \gamma, \delta, \varepsilon
\]
one obtains the following table:
\[
\begin{array}{c!{\vrule width 0.8pt}c|c|c|c|c|c}
 & 1 & \alpha & \beta & \gamma & \delta & \varepsilon \\
\noalign{\hrule height 0.8pt}
1 & 1 & \alpha & \beta & \gamma & \delta & \varepsilon \\ \hline
\alpha & \alpha & \beta & 1 & \varepsilon & \gamma & \delta \\ \hline
\beta & \beta & 1 & \alpha & \delta & \varepsilon & \gamma \\ \hline
\gamma & \gamma & \delta & \varepsilon & 1 & \alpha & \beta \\ \hline
\delta & \delta & \varepsilon & \gamma & \beta & 1 & \alpha \\ \hline
\varepsilon & \varepsilon & \gamma & \delta & \alpha & \beta & 1
\end{array}
\]

In every row and in every column, as must be the case, all the elements occur.

\begin{center}
\S.\,28.

\gesperrt{Relation of the general groups to the permutation groups.}
\end{center}

The groups that we have hitherto employed most are the permutation groups; this kind of groups acquires a heightened significance for general group theory as well through the following considerations.

Let
\weq{1}{P = a_0, a_1, a_2, \ldots, a_{n-1}}
be an arbitrary group of degree $n$. If we take from $P$ any
%CONT>>
